\section{Related Work}\label{sec:related}
\noindent\textbf{Multi-Agent Coordination and Specialization.} The classic characterization of agency by Wooldridge and Jennings defines agents through autonomy, reactivity, pro-activeness, and social ability~\cite{wj95}. In modern LLM-based systems, multi-agent frameworks predominantly follow two coordination paradigms: conversational emergent interaction, exemplified by AutoGen~\cite{autogen}, and structured role-based workflows with standardized operating procedures, exemplified by MetaGPT~\cite{metagpt}. While these frameworks demonstrate qualitative versatility across software development and debate tasks, recent empirical audits demonstrate that multi-agent LLM systems frequently suffer from coordination collapse, specification drift, and inter-agent misalignment~\cite{mast,resil,masfire}. Crucially, literature often assumes that decomposing a complex task into multiple specialized agents is inherently advantageous. However, few studies isolate whether specialization provides measurable decision-making benefits over a unified controller operating with equivalent information. Our work directly evaluates this trade-off by comparing specialized stateful agents against an uncoordinated baseline and an equivalent unified architecture.

\noindent\textbf{Adaptive Mock Interview and Tutoring Systems.} AI-assisted interview systems have emerged to provide automated feedback and dynamic evaluation for technical candidates. Systems such as Zara~\cite{zara} utilize generative LLMs to synthesize candidate feedback, while empirical user studies investigate student confidence, readiness, and trust in automated code evaluation~\cite{virt,chase}. However, generative LLMs often exhibit verbosity, position bias, and hallucinations when tasked with simultaneous scoring and pedagogical planning~\cite{judge}. To ensure evaluation integrity, our target system enforces strict authority boundaries derived from the principle of least privilege~\cite{ss75}: scoring is computed by a deterministic, cross-encoder evaluator outside the language model, while agents operate solely on validated observations.

\noindent\textbf{Reinforcement Learning, Shielding, and Coordination Envelopes.} Dynamic difficulty adaptation in educational platforms is commonly framed as a reinforcement learning problem. Single-agent RL approaches, such as Proximal Policy Optimization (PPO)~\cite{ppo} implemented via Stable-Baselines3~\cite{sb3}, optimize policy returns over state trajectories. However, unconstrained RL policies risk erratic behavior during out-of-distribution candidate interactions. Safe reinforcement learning surveys~\cite{safe} classify external knowledge mechanisms, such as shielding~\cite{shield}, as essential safety envelopes that override learned actions when constraints are violated. Our arbitration architecture implements an empirical safety envelope where rule-based, evidenced agent proposals compete with the frozen PPO action while remaining bounded by deterministic guardrails. This architecture represents centrally orchestrated role specialization rather than Multi-Agent Reinforcement Learning (MARL)~\cite{marl}, as the PPO policy is frozen and the agents do not update policy weights dynamically.

\noindent\textbf{Methodological Foundations and Dependability.} To evaluate coordination integrity, we draw upon dependable computing taxonomies~\cite{alrl04}, enforcing bulkhead isolation via per-agent circuit breakers~\cite{nygard} to bound cascading failures. For statistical inference over replayed candidate answer clusters, we employ clustered bootstrapping methodologies established by Field and Welsh~\cite{fw07}, preventing artificially narrow confidence intervals that arise from treating permuted question sequences as independent candidate samples.
